% =====================================================================
%  CLASE 12 -- Sistemas de drenaje de aguas lluvias
% =====================================================================
\clase{12}{Sistemas de drenaje urbano de aguas lluvias}

\section{Motivación}
Las inundaciones de calles en temporales son frecuentes en ciudades con
sistemas de drenaje insuficientes o mal mantenidos (sumideros tapados, calles
anegadas). Ver: \emph{La historia del alcantarillado del Santiago del
Centenario y su resistencia ante la lluvia}.

\section{Sistemas de drenaje de aguas lluvias}
El objetivo de estos sistemas es \textbf{evitar la inundación de las calles}:
el agua escurre por las calles (con pendientes $S_x$, $S_y$), es captada por
sumideros y conducida por colectores hasta un cauce receptor (río).

\begin{figure}[H]
\centering
\begin{tikzpicture}[font=\small]
  \fill[gray!30,rounded corners=8pt] (-0.3,-0.3) rectangle (6.9,6.9);
  \foreach \i in {0,1,2} \foreach \j in {0,1,2}
    \fill[suelo!80,rounded corners=4pt] ({0.2+2.2*\i},{0.6+2.2*\j}) rectangle ({1.8+2.2*\i},{2.2+2.2*\j});
  \foreach \i in {0,1,2} \foreach \j in {0,1,2}{
    \draw[fill=white] ({2.0+2.2*\i},{0.4+2.2*\j}) circle (0.17);
    \draw[flecha,azulusm] ({0.3+2.2*\i},{0.4+2.2*\j}) -- ({1.75+2.2*\i},{0.4+2.2*\j});
    \draw[flecha,azulusm] ({2.0+2.2*\i},{2.1+2.2*\j}) -- ({2.0+2.2*\i},{0.65+2.2*\j});}
  \draw[flecha,azulusm,line width=2pt] (6.4,0.4) -- (7.3,-0.4) node[right]{Río};
  \draw[flecha] (3,7.1) -- (4,7.1) node[midway,above]{$S_x$};
  \draw[flecha] (-0.5,5.5) -- (-0.5,4.5) node[midway,left]{$S_y$};
\end{tikzpicture}
\caption{Esquema de una red urbana: escurrimiento por las calles, captación
en sumideros/cámaras (círculos) y descarga al río.}
\end{figure}

\begin{itemize}
  \item Las obras de drenaje están diseñadas para captar, retener, almacenar,
        infiltrar, transportar y descargar aguas lluvias.
  \item Las obras hidráulicas de drenaje son para redes domiciliaria,
        secundaria y primaria.
  \item En la red domiciliaria y secundaria se da mayor importancia a la
        infiltración y al almacenamiento.
  \item En la red primaria se da relevancia al transporte y descarga en las
        redes de aguas abajo.
\end{itemize}
Referencia: \emph{Manual de Drenaje Urbano (MDU)}, DOH: Cap.~4 Estudios
Básicos; Cap.~5 Técnicas de Diseño de Redes; Cap.~6 Diseño Hidráulico de Obras.

\begin{figure}[H]
\centering
\begin{tikzpicture}[font=\scriptsize,node distance=0.5cm,
  caja/.style={draw,fill=#1,minimum width=1.7cm,minimum height=0.7cm,align=center}]
  \node[caja=gray!20] (r) {Redes};
  \node[caja=rojo!20,right=of r] (d) {Red\\domiciliaria};
  \node[caja=suelo!40,right=of d] (s) {Red\\secundaria};
  \node[caja=green!25,right=of s] (p) {Red\\primaria};
  \node[caja=aguaclara,right=of p] (n) {Red\\natural};
  \node[caja=gray!20,below=of r] (e) {Elementos\\propios};
  \node[caja=rojo!20,below=of d] (d2) {Techos\\pavimentos};
  \node[caja=suelo!40,below=of s] (s2) {Calles\\sumideros};
  \node[caja=green!25,below=of p] (p2) {Cauces\\descargas};
  \node[caja=aguaclara,below=of n] (n2) {Zonas\\inundables};
  \draw[->] (d) -- (s); \draw[->] (s) -- (p); \draw[->] (p) -- (n);
  \foreach \a/\b in {d2/d,s2/s,p2/p,n2/n} \draw[->] (\a) -- (\b);
\end{tikzpicture}
\caption{Organización de las redes de drenaje urbano (MDU, Fig.~5.1.2).
Complementos: infiltración, almacenamiento y conducción.}
\end{figure}

Estos sistemas están compuestos por: \textbf{cunetas}, \textbf{sumideros},
\textbf{cámaras de observación} y \textbf{colectores de aguas lluvias}.

\begin{figure}[H]
\centering
\begin{tikzpicture}[font=\small]
  \fill[gray!50] (0,0) -- (0,1.8) -- (0.3,1.8) -- (0.3,0.9) -- (4,1.8) -- (7.7,0.9) -- (7.7,1.8) -- (8,1.8) -- (8,0) -- cycle;
  \fill[aguaclara] (0.3,0.9) -- (0.3,1.25) -- (1.7,1.25) -- cycle;
  \fill[aguaclara] (7.7,0.9) -- (7.7,1.25) -- (6.3,1.25) -- cycle;
  \fill[sueloclaro] (0.5,-1.3) rectangle (7.5,0);
  \draw[thick,fill=white] (4,-0.6) circle (0.5);
  \begin{scope}\clip (4,-0.6) circle (0.5); \fill[aguaclara] (3.4,-1.2) rectangle (4.6,-0.75);\end{scope}
  \draw[thick] (4,-0.6) circle (0.5);
  \draw (1.0,0.9) -- (1.8,0.9); \draw (1.8,0.9) arc (0:14:0.8); \node at (2.1,1.0) {$i$};
  \draw[<-] (0.7,1.1) -- (-0.8,2.2) node[above,azulusm,align=center]{Pendiente transversal\\o bombeo};
  \draw[<-] (4.5,-0.8) -- (8.2,-1.0) node[right,azulusm,align=left]{Colector de\\aguas lluvias};
  \draw[->,azulusm,thick] (1.2,1.0) .. controls (2.2,0.2) .. (3.5,-0.5);
  \draw[->,azulusm,thick] (6.8,1.0) .. controls (5.8,0.2) .. (4.5,-0.5);
\end{tikzpicture}
\caption{Perfil transversal de una calle: cunetas a ambos lados (bombeo $i$)
que descargan, vía sumideros, al colector.}
\end{figure}

\section{Tipos de obras}
Las obras dependen de la capacidad de conducción ($Q$) requerida:
\begin{itemize}
  \item Cuneta simple.
  \item Cuneta compuesta.
  \item Colector de aguas lluvias (con sumideros y cámaras).
\end{itemize}

% Macro sección de cuneta: #1=1 compuesta
\newcommand{\seccioncuneta}[1]{%
\begin{tikzpicture}[font=\small]
  \fill[gray!55] (0,0) rectangle (0.3,1.7);
  \ifnum#1=1
    \fill[gray!40] (0.3,0) -- (0.3,0.1) -- (1.8,0.9) -- (5.5,1.4) -- (6.3,1.25) -- (6.3,0) -- cycle;
    \fill[aguaclara] (0.3,0.1) -- (0.3,1.1) -- (3.5,1.1) -- (1.8,0.9) -- cycle;
    \draw (0.3,0.1) -- (1.8,0.9) -- (5.5,1.4) -- (6.3,1.25);
    \draw[dashed] (1.8,0.9) -- (5.8,0.9) (1.8,-0.2) -- (1.8,0.9);
    \draw[cota] (0.3,-0.1) -- (1.8,-0.1) node[midway,below]{$w$};
    \node at (1.3,0.3) {$i_1$}; \node at (4.9,1.1) {$i_2$};
    \draw[cota] (0.3,1.35) -- (3.5,1.35) node[midway,above]{$b$};
  \else
    \fill[gray!40] (0.3,0) -- (0.3,0.4) -- (5.5,1.2) -- (6.3,1.05) -- (6.3,0) -- cycle;
    \fill[aguaclara] (0.3,0.4) -- (0.3,0.85) -- (3.2,0.85) -- cycle;
    \draw (0.3,0.4) -- (5.5,1.2) -- (6.3,1.05);
    \draw[dashed] (0.3,0.4) -- (5.8,0.4);
    \node at (4.9,0.6) {$i$};
    \draw[cota] (0.3,1.1) -- (3.2,1.1) node[midway,above]{$b$};
    \draw[cota] (0.1,0.4) -- (0.1,0.85) node[midway,left]{$y$};
    \draw[cota] (-0.4,0.4) -- (-0.4,1.7) node[midway,left]{$h_s$};
  \fi
  \draw[<-] (5.5,1.45) -- (5.9,1.9) node[above,azulusm]{Bombeo};
\end{tikzpicture}}

\section{Cuneta simple}
Conformada por la calle misma y la solera.
\begin{itemize}
  \item Destinada para caudales pequeños.
  \item Para su diseño se debe determinar la pendiente transversal de la
        calle ($i$, bombeo) en función de:
  \begin{itemize}
    \item Ancho máximo permitido ($b_{max}=1-2$ [m]).
    \item Pendiente longitudinal de la calle ($S=3\%-10\%$).
    \item Revisar Tabla 5.3.4 del MDU.
  \end{itemize}
\end{itemize}

\begin{figure}[H]
\centering
\seccioncuneta{0}
\caption{Corte transversal de una cuneta simple: $y=b\cdot i\leq h_s\approx15$ [cm].}
\end{figure}

Utilizando la ecuación de Manning:
\[
Q=\frac{\sqrt S}{n}\Omega\cdot R_H^{2/3},\qquad
\Omega=\frac{b\cdot y}{2}=\frac{b^2\cdot i}{2},\qquad R_H=\frac\Omega\Psi,\qquad
\Psi=b\cdot i+b\sqrt{i^2+1}\approx b(i+1)
\]
\begin{formula}
$\displaystyle Q=\frac{\sqrt S}{n}\,\frac{b^2\cdot i}{2}\left(\frac{b\cdot i}{2(i+1)}\right)^{2/3}$\qquad $i$ varía entre 2 -- 4\%
\end{formula}

\section{Cuneta compuesta}
Se usan cuando el ancho de inundación está en el rango de 1 a 2 [m] para una
cuneta simple y se requiere mayor capacidad.
\begin{itemize}
  \item Generalmente utilizadas en carreteras.
  \item El ancho superficial ($b$) se calcula con la fórmula de Manning,
        teniendo en cuenta que $b\leq1-2$ [m].
  \item Revisar Tabla 5.3.5 del MDU.
  \item Si $b<w$, el cálculo se debe realizar como cuneta simple.
\end{itemize}

\begin{figure}[H]
\centering
\seccioncuneta{1}
\caption{Corte transversal de una cuneta compuesta (con zarpa de ancho $w$).}
\end{figure}

\begin{formula}
$\displaystyle Q=0.315\frac{\sqrt S}{n}b^{8/3}\left[i_2+(i_1-i_2)\left(\frac wb\right)^2\right]^{5/3}$
\end{formula}
con $w$: ancho de la zarpa ($w\approx0.6$ [m]); $S$: pendiente longitudinal
de la calle; $i_2$: bombeo de la calle (2 -- 4\%); $i_1$: pendiente
transversal de la zarpa (2 -- 10\%).

\section{Sumideros}
Obras de arte que se utilizan para captar el agua de la calle hacia el colector.
\begin{table}[H]
\centering\small
\begin{tabularx}{\textwidth}{lX}
\toprule
\textbf{Tipo} & \textbf{Características}\\
\midrule
Horizontal de fondo (tipo S3 y S4, SERVIU) & Usado para un amplio rango de pendientes. Inconvenientes: peatones, ciclistas, obstrucción.\\
Lateral (tipo S2, SERVIU) & Funciona admitiendo basuras; su capacidad decrece para pendientes longitudinales mayores al 3\%, donde se recomienda cuneta con pendiente longitudinal para mayor captación.\\
Mixto (tipo S1 y S2, SERVIU) & Funciona admitiendo basuras; capacidad de captación aumenta 10\% con respecto al sumidero horizontal.\\
\bottomrule
\end{tabularx}
\end{table}

\begin{figure}[H]
\centering
\begin{tikzpicture}[font=\small,x={(1cm,0cm)},y={(0.5cm,0.35cm)},z={(0cm,1cm)}]
  % losa
  \draw[fill=gray!15] (0,0,0) -- (5,0,0) -- (5,4,0) -- (0,4,0) -- cycle;
  \draw[fill=gray!35] (0,4,0) -- (5,4,0) -- (5,4,0.6) -- (0,4,0.6) -- cycle;
  % rejilla
  \draw[fill=gray!60] (1.2,2.4,0) -- (3.8,2.4,0) -- (3.8,3.6,0) -- (1.2,3.6,0) -- cycle;
  \foreach \x in {1.45,1.7,...,3.6} \draw[white] (\x,2.45,0) -- (\x,3.55,0);
  \draw[cota] (1.2,2.1,0) -- (3.8,2.1,0) node[midway,below]{$L$};
  \draw[cota] (4.1,2.4,0) -- (4.1,3.6,0) node[midway,right]{$w$};
\end{tikzpicture}
\caption{Sumidero horizontal de fondo junto a la solera: reja de largo $L$ y
ancho $w$ ($w=0.4\sim0.7$ [m], $L=1\sim2$ [m]).}
\end{figure}

\subsection{Sumidero horizontal de fondo}
El caudal máximo de captación depende del funcionamiento:
\begin{enumerate}
  \item \textbf{Actuando como vertedero} (alturas de agua pequeñas): un
        vertedero es un punto de control; produce un peralte del
        escurrimiento aguas arriba.
  \begin{formula}
  $\displaystyle Q_m=1.66\,(L_e+2\cdot w_e)\,h^{1.5}\ [\text{m}^3/\text{s}]\qquad\text{si } h<1.6\frac{A_e}{L_e+2w_e}$
  \end{formula}
  \item \textbf{Actuando como orificio} (alturas de agua mayores): el
        sumidero se ahoga.
  \begin{formula}
  $\displaystyle Q_m=2.66\cdot A_e\cdot h^{0.5}\ [\text{m}^3/\text{s}]\qquad\text{si } h\geq1.6\frac{A_e}{L_e+2w_e}$
  \end{formula}
\end{enumerate}
donde $w_e=w-e\cdot n_L$: ancho efectivo [m]; $L_e=L-e\cdot n_T$: largo
efectivo [m]; $A_e=L_e\cdot w_e$: área efectiva [m$^2$]; $n_L$: número de
barras longitudinales; $n_T$: número de barras transversales; $e$: diámetro o
espesor de las barras [m]; $h$: altura de agua en la calle frente al sumidero [m].

\paragraph{Eficiencia.} El sumidero no siempre logra captar toda el agua que
circula por la cuneta:
\begin{formula}
$Q_s=\begin{cases}\eta_HQ & \eta_HQ<Q_m\\ Q_m & \eta_HQ\geq Q_m\end{cases}$
\end{formula}
\[
\eta_H=R_f\cdot\eta_0+R_s(1-\eta_0),\qquad
\eta_0=1-\left(1-\frac w{b_s}\right)^{2.67},
\]
\[
R_s=\frac{1}{1+\dfrac{0.0828\,v^{1.8}}{i\cdot L^{2.3}}},\qquad
R_f=1-0.295(v-v_0)
\]
todos acotados entre 0 y 1; $\eta_H$: eficiencia del sumidero de fondo;
$Q_s$: caudal captado [m$^3$/s]; $Q$: caudal que circula por la cuneta
[m$^3$/s]; $v$: velocidad del agua en la calle frente al sumidero [m/s];
$v_0$: velocidad a la que ocurre el primer fenómeno de salpicadura en la reja [m/s].

\subsection{Sumidero lateral}
\begin{formula}
$Q_m=\begin{cases}1.27\cdot L\cdot h^{1.5} & h<a\quad\text{(vertedero)}\\ 2.66\cdot L\cdot a\cdot h^{0.5} & h\geq a\quad\text{(ahogado)}\end{cases}$
\end{formula}
Recomendable su uso sólo para pendientes longitudinales $S<3\%$.

\paragraph{Eficiencia}
\begin{formula}
$Q_s=\begin{cases}\eta_LQ & \eta_LQ<Q_m\\ Q_m & \eta_LQ\geq Q_m\end{cases}\qquad
\eta_L=\begin{cases}1 & h>a\\ 1-\left(1-\dfrac{L}{L_T}\right)^{1.8} & h\leq a\end{cases}$
\end{formula}
\[
\begin{aligned}
&\text{Cuneta compuesta:} && L_T=0.81\,Q^{0.42}\cdot S^{0.3}(n\cdot i)^{-0.6}\ [\text{m}]\\
&\text{Cuneta simple } (w=0): && L_T=0.81\,Q^{0.42}\cdot S^{0.3}(n\cdot i_2)^{-0.6}\ [\text{m}]
\end{aligned}
\]
con $n$: rugosidad de Manning; $S$: pendiente longitudinal de la calle;
$L_T$: largo necesario para captar todo el caudal ($L_{T,min}=L$).

\begin{figure}[H]
\centering
\begin{tikzpicture}[font=\small]
  \fill[gray!30] (0,0) rectangle (2.5,5);
  \fill[gray!55] (-0.2,0) rectangle (0,5); \fill[gray!55] (2.5,0) rectangle (2.7,1.6); \fill[gray!55] (2.5,3.0) rectangle (2.7,5);
  \draw[dashed,thick] (1.2,0) -- (1.2,5);
  \draw[flecha,azulusm] (2.0,4.9) -- (2.0,4.0) node[pos=0,right]{$Q$};
  \foreach \y in {3.6,3.1,2.6,2.1} \draw[->,azulusm] (2.0,\y+0.3) .. controls (2.3,\y) .. (2.5,\y-0.1);
  \draw[flecha,azulusm,line width=2pt] (2.5,2.3) -- (3.3,2.3) node[right]{$\eta_LQ$};
  \draw[cota] (2.9,1.6) -- (2.9,3.0) node[pos=0.8,right]{$L$};
  \draw[cota] (4.2,1.2) -- (4.2,3.6) node[midway,right]{$L_T$};
  \draw[dashed] (2.7,1.2) -- (4.3,1.2) (2.7,3.6) -- (4.3,3.6);
  \node[below] at (1.8,0) {$(1-\eta_L)Q$};
\end{tikzpicture}
\caption{Sumidero lateral (vista en planta): capta $\eta_LQ$ y deja pasar $(1-\eta_L)Q$.}
\end{figure}

\subsection{Sumideros mixtos}
\begin{itemize}
  \item Combinación de ambos tipos (de fondo y lateral).
  \item La experiencia indica que sólo existe un pequeño aumento del caudal
        de captación respecto a sumideros de fondo ($\approx10\%$).
  \item Se recomienda su uso en zonas con muchos árboles de hoja caduca, ya
        que de emergencia puede funcionar el sumidero lateral.
\end{itemize}
\begin{formula}
$Q_s=\begin{cases}(\eta_H+\eta_L)Q & (\eta_H+\eta_L)Q<Q_m\\ Q_m & (\eta_H+\eta_L)Q\geq Q_m\end{cases}$
\end{formula}

\subsection{Factor de seguridad de eficiencia}
Las eficiencias empíricas o calculadas suponen que el sumidero está limpio.
En la práctica esta situación rara vez ocurre; por lo tanto, es importante
incluir un factor de seguridad $F$:
\begin{formula}
$\eta_d=F\cdot\eta$
\end{formula}
\begin{table}[H]
\centering
\begin{tabular}{lcc}
\toprule
\textbf{Caudal aportante} & \textbf{Mantención y limpieza} & \textbf{Factor de seguridad}\\
\midrule
Inferior a 30 l/s & Buena & 1.0\\
 & Mala & 0.8\\
Superior a 30 l/s & Buena & 0.7\\
 & Mala & 0.6\\
\bottomrule
\end{tabular}
\caption{Estimación de un factor de seguridad para sumideros (MDU, Tabla 6.5.13).}
\end{table}
Caudales altos tienden a llevar más basura y por lo tanto a taparse; por ello
se recomienda que siempre se haga con ventana lateral, así ésta no se tenga
en cuenta para el cálculo de la eficiencia.

\subsection{Distancia entre sumideros}
\begin{itemize}
  \item Se debe garantizar que el flujo no pase las condiciones de seguridad.
  \item Para la red primaria, el caudal máximo en la calle debe calcularse
        mediante el modelo lluvia--escorrentía SWMM.
  \item Para la red secundaria, el flujo que llega a la calle se puede
        estimar por el método racional.
\end{itemize}
\begin{formula}
$\displaystyle Q_{captado}=\frac{C\cdot i\cdot A}{3600}=\eta_dQ_{max}
\quad\Longrightarrow\quad
L=\frac{A}{b_c}=\frac{3600\,(F\cdot\eta)\,Q_{max}}{C\cdot i\cdot b_c}$
\end{formula}
\[
\text{Calzada con doble bombeo: } b_c=0.5\cdot e+w_c;\qquad\text{calle con bombeo único: } b_c=e+w_c
\]
donde $L$: separación entre sumideros [m]; $A$: área aportante [m$^2$];
$w_c$: ancho de los terrenos aportantes [m]; $e$: ancho de la calle [m]; $C$:
coeficiente de escorrentía; $i$: intensidad de diseño [mm/h] ($T_r=2$ años,
$D_{min}=30$ minutos).

\subsection{Consideraciones de diseño (cunetas y sumideros)}
\begin{itemize}
  \item Para tormentas menores ($T_r=2$ a 10 años), la calzada no podrá
        inundarse más del ancho recomendado. Para el cálculo del tiempo de
        concentración se pueden utilizar las fórmulas vistas en hidrología,
        pero se debe verificar $t_c\geq30$ [min].
  \item Para tormentas mayores ($T_r=50$ a 100 años), la calzada no podrá
        inundarse con altura mayor a la solera y en todos los casos $h<15$ cm.
\end{itemize}

\section{Colectores de aguas lluvias}
\begin{itemize}
  \item Se utilizan en el caso de que $b>1-2$ [m] (excede lo permitido para una cuneta).
  \item Su diseño considera caudales ($Q$) pequeños y medianos.
  \item Se utilizan tuberías de diámetro mínimo de 300 [mm] y máximo 2 [m]:
        tubos o cajones de cemento comprimido (C.C.C.) o tubos de HDPE.
  \item MDU 6.7.2: colectores de diámetro pequeño ($D<1.2$ m), $T_r=2-10$ años.
  \item MDU 6.7.2: colectores de diámetro mayor ($D>1.2$ m), $T_r=5-10$ años.
\end{itemize}

La elección del tamaño queda definida por el caudal y los requisitos de
velocidad del flujo. La altura se puede calcular con la fórmula de Manning:
\[
Q=\frac{\sqrt S}{n}A\cdot R_H^{2/3}
\]
\begin{itemize}
  \item Bajo condiciones normales se debe verificar que $h\leq0.8D$.
  \item Para el diseño se considera la condición crítica ($h=0.8D$).
  \item Bajo condiciones excepcionales podría exceder el límite de $0.8D$
        (p.~ej.\ $T_r=50$ años), pero siempre sin presión.
\end{itemize}

\begin{figure}[H]
\centering
\begin{minipage}{0.35\textwidth}\centering
\begin{tikzpicture}[font=\small]
  \begin{scope}\clip (0,0) circle (1.2); \fill[aguaclara] (-1.3,-1.3) rectangle (1.3,0.72);\end{scope}
  \draw[thick] (0,0) circle (1.2);
  \draw (-0.96,0.72) -- (0.96,0.72);
  \draw (-0.96,0.72) -- (0,0) -- (0.96,0.72);
  \draw[->] (0,0) ++(-127:0.4) arc (-127:127:0.4) node[pos=0.5,right]{$\theta$};
  \pic at (0,0.74) {nivelagua};
  \draw[cota] (-1.6,-1.2) -- (-1.6,1.2) node[midway,left]{$D$};
  \draw[cota] (1.5,-1.2) -- (1.5,0.72) node[midway,right]{$h$};
\end{tikzpicture}
\end{minipage}
\begin{minipage}{0.55\textwidth}
\[
\begin{gathered}
A=\frac{(\theta-\operatorname{sen}\theta)\,D^2}{8},\qquad
R_h=\left(1-\frac{\operatorname{sen}\theta}{\theta}\right)\frac D4\\
\theta=2\cos^{-1}\left(1-\frac{2h}{D}\right)
\end{gathered}
\]
\end{minipage}
\caption{Sección circular parcialmente llena.}
\end{figure}

\begin{table}[H]
\centering
\begin{tabular}{lccccc}
\toprule
\textbf{Situación} & $h$ [m] & $\theta$ [grados] & $A_m$ [m$^2$] & $P_m$ [m] & $R_h$ [m]\\
\midrule
Sección completa & $D$ & 360 & $0.785D^2$ & $3.142D$ & $0.250D$\\
Caudal máximo & $0.95D$ & 308 & $0.771D^2$ & $2.686D$ & $0.287D$\\
Sección diseño & $0.80D$ & 254 & $0.674D^2$ & $2.217D$ & $0.304D$\\
\bottomrule
\end{tabular}
\caption{Propiedades de secciones especiales en ductos circulares (MDU, Tabla 6.7.13).}
\end{table}

Además, se debe verificar:
\begin{itemize}
  \item \textbf{Velocidad mínima} (evitar depósito de sedimentos dentro de la
        tubería, autolavado):
  \begin{formula}
  $\displaystyle V_c=0.397\,\frac{D^{2/3}\sqrt S}{n}\geq0.6\ [\text{m/s}]$
  \end{formula}
  \item \textbf{Velocidad máxima} (evitar socavación):
\end{itemize}
\begin{table}[H]
\centering
\begin{tabular}{lc}
\toprule
\textbf{Tipo de tubería} & \textbf{Velocidad máxima} [m/s]\\
\midrule
Concreto simple hasta 45 cm de diámetro & 3.0\\
Concreto reforzado de 61 cm de diámetro o mayores & 3.5\\
Fibrocemento & 5.0\\
Policloruro de vinilo (PVC) & 5.0\\
Polietileno de alta densidad & 5.0\\
Tubería de acero corrugado & 4.5\\
Tubería de acero corrugado revestido en shotcrete & 3.0\\
Polímero reforzado con fibra de vidrio (PRFV) & 3.5\\
\bottomrule
\end{tabular}
\caption{Velocidad máxima permisible en colectores de pequeño diámetro (MDU,
Tabla 6.7.14; fuente: Comisión Nacional del Agua, México, 2009).}
\end{table}

\section{Recomendaciones de diseño de colectores}
\paragraph{Pendientes longitudinales muy grandes.} Si se tienen pendientes
$S$ muy grandes, dentro de las soluciones para reducir la velocidad se encuentran:
\begin{itemize}
  \item \textbf{Cámaras de observación:} se colocan para escalonar el
        escurrimiento, de forma tal que las tuberías no registren velocidades
        mayores a lo permitido ($v_{max}$), con $S_D<S$.
  \item \textbf{Dientes disipadores de energía:} reducen la altura de
        velocidad; se usan en cajones de hormigón; ¡alto costo!
\end{itemize}

\begin{figure}[H]
\centering
\begin{tikzpicture}[font=\small]
  \fill[gray!50] (0,3.0) -- (9,1.2) -- (9,1.4) -- (0,3.2) -- cycle;
  \node[above] at (4.5,2.3) {$S$};
  \foreach \x/\yt in {2/2.6,5.5/1.9}{
    \fill[gray!25,draw] (\x,\yt) rectangle (\x+0.6,\yt-2.0);}
  \draw[thick] (0,1.6) -- (2,1.25) (0,1.8) -- (2,1.45);
  \draw[thick] (2.6,0.8) -- (5.5,0.35) (2.6,1.0) -- (5.5,0.55);
  \draw[thick] (6.1,-0.1) -- (9,-0.5) (6.1,0.1) -- (9,-0.3);
  \node at (1,1.2) {$S_D$}; \node at (4,0.3) {$S_D$}; \node at (7.5,-0.6) {$S_D$};
  \node[azulusm,align=center] at (8,2.4) {Cámaras de\\observación};
  \draw[->,azulusm] (7.2,2.3) -- (6.1,1.4);
  \node at (4,-0.8) {$S_D<S$};
\end{tikzpicture}
\caption{Cámaras de observación para escalonar un colector en una calle de gran pendiente.}
\end{figure}

\paragraph{Trampas de arena.} Se deben colocar si existen:
\begin{itemize}
  \item Cambios de dirección.
  \item Cambios de pendiente.
  \item Tramos largos sin ninguna perturbación en el flujo (cámaras cada
        $50\sim130$ [m]).
\end{itemize}
La trampa de arena es la zona de la cámara ubicada bajo la profundidad de la
tubería; retiene los sedimentos que trae el flujo y se debe limpiar cada
cierto tiempo.

\paragraph{Profundidad de instalación.} Colector de aguas lluvias a 1.8 [m]
bajo la superficie; red de agua potable a $1\sim1.5$ [m].

\paragraph{Diámetros mínimos.} Para posibilitar la limpieza y correcto
funcionamiento. Desde el punto de vista administrativo, existen dos tipos de
colectores:
\begin{itemize}
  \item Primarios ($D>0.8$ [m]): dependen de la Dirección de Obras Hidráulicas (DOH).
  \item Secundarios ($D\leq0.8$ [m]): dependen del SERVIU.
\end{itemize}

\section{Obtención de caudales de diseño}
Se tiene una serie de cuencas conectadas a través de una red de colectores.
Como las cuencas son de tamaño pequeño, se asumen condiciones meteorológicas
homogéneas, es decir, la precipitación es la misma en todas las cuencas.
Generalmente se utiliza la fórmula racional para el cálculo de $Q$.

\begin{figure}[H]
\centering
\begin{tikzpicture}[font=\small]
  \foreach \x/\y/\n in {0/2.2/1,3.5/1.8/2,7/1.5/3}{
    \fill[green!55!black!60,draw=green!30!black] (\x,\y) .. controls (\x-0.9,\y+0.6) and (\x-0.7,\y+1.8) .. (\x+0.1,\y+1.9)
      .. controls (\x+0.9,\y+2.0) and (\x+1.1,\y+0.9) .. (\x,\y);
    \node at (\x+1.0,\y+1.7) {$A_\n$};
    \fill[white,draw] (\x,\y) circle (0.07);}
  \draw[flecha,azulusm,thick] (0,2.2) -- (3.5,1.8) node[midway,above]{$Q_{D1},L_1$} node[midway,below=4pt,draw,circle,inner sep=1pt]{1};
  \draw[flecha,azulusm,thick] (3.5,1.8) -- (7,1.5) node[midway,above]{$Q_{D2},L_2$} node[midway,below=4pt,draw,circle,inner sep=1pt]{2};
  \draw[flecha,azulusm,thick] (7,1.5) -- (9.5,1.2) node[midway,above]{$Q_{D3},L_3$} node[midway,below=4pt,draw,circle,inner sep=1pt]{3} node[right]{Río};
\end{tikzpicture}
\caption{Serie de cuencas conectadas por colectores.}
\end{figure}

\paragraph{Colector 1:} sólo aporta caudal la cuenca 1:
\[
Q_1=\frac{C_1\cdot\bar\imath(t_c)\cdot A_1}{3.6}\quad\Longrightarrow\quad Q_{D1}=Q_1
\]

\paragraph{Colector 2:} las cuencas 1 y 2 aportan caudal al colector 2. El
caudal de diseño corresponde a:
\begin{formula}
$Q_{D2}=\max(Q_{D1},Q_2,Q_{2I})$
\end{formula}
donde $Q_{D1}$: caudal de diseño del colector 1; $Q_2$: caudal máximo
aportado por la cuenca 2; $Q_{2I}$: caudal máximo aportado por las cuencas 1 y 2:
\[
Q_{2I}=\frac{\bar C_{1,2}\cdot\bar\imath(t_{c2I})\cdot(A_1+A_2)}{3.6},\qquad
\bar C_{1,2}=\frac{C_1+C_2}{2},\qquad
t_{c2I}=\max(t_{c1}+t_{t1},\,t_{c2}),\qquad t_{t1}=\frac{L_1}{v_1}
\]

\paragraph{Colector 3:} el procedimiento es similar al caso del colector 2:
\begin{formula}
$Q_{D3}=\max(Q_{D2},Q_3,Q_{3I})$
\end{formula}
\[
Q_{3I}=\frac{\bar C_{1,2,3}\cdot\bar\imath(t_{c3I})\cdot(A_1+A_2+A_3)}{3.6},\qquad
\bar C_{1,2,3}=\frac{C_1+C_2+C_3}{3},
\]
\[
t_{c3I}=\max(t_{c2I}+t_{t2},\,t_{c3}),\qquad t_{t2}=\frac{L_2}{v_2}
\]

\paragraph{Hidrograma unitario triangular del SCS.} Si se dispone de la
distribución temporal de la precipitación (hietograma), una alternativa es
utilizar el hidrograma de crecidas del SCS:
\begin{formula}
$\displaystyle t_p=\frac{t_{LL}}{2}+0.6\cdot t_c\ \ [\text{h}],\qquad t_r=1.67\cdot t_p\ \ [\text{h}],\qquad t_B=2.67\cdot t_p\ \ [\text{h}]$
\end{formula}
Igualando el volumen de escorrentía directa con el área del hidrograma:
\[
V_{ed}=P_{ef}\cdot A=\frac12Q_p\cdot t_B
\quad\Longrightarrow\quad
\boxed{Q_p=\frac{2\cdot P_{ef}\cdot A}{2.67\left(\frac{t_{LL}}{2}+0.6\cdot t_c\right)}}
\]

\begin{figure}[H]
\centering
\begin{tikzpicture}[font=\small]
  \draw[flecha] (0,0) -- (7,0) node[right]{$t$};
  \draw[flecha] (0,0) -- (0,3) node[left]{$Q$};
  \fill[azulusm] (0,3.4) rectangle (0.6,3.2); \node[right] at (0.6,3.3) {$P_{ef}$ (duración $t_{LL}$)};
  \draw[azulusm,thick] (0,0) -- (2.2,2.4) -- (6,0);
  \draw[dashed] (0,2.4) node[left]{$Q_p$} -- (2.2,2.4) -- (2.2,0);
  \draw[cota] (0,-0.3) -- (2.2,-0.3) node[midway,below]{$t_p$};
  \draw[cota] (2.2,-0.3) -- (6,-0.3) node[midway,below]{$t_r$};
  \draw[cota] (0,-0.9) -- (6,-0.9) node[midway,below]{$t_B$};
\end{tikzpicture}
\caption{Hidrograma unitario triangular del SCS.}
\end{figure}
